% SPDX-License-Identifier: LPPL-1.3c
% Copyright (C) 2026 Christophe Jorssen
% Author and Current Maintainer: Christophe Jorssen <christophe.jorssen@gmail.com>
% LPPL maintenance status: maintained
% This file is part of luafemm. See LICENSE and LICENSES.md for its terms.
\section{Lua API and development}
\label{sec:development}
\subsection{Module boundaries}
Every Lua file returns a module table and uses local implementation
variables; there is no deprecated |module()| call and no runtime change
to |package.path|. LuaTeX's or ConTeXt's TEXMF resolver finds installed
files. The TeX loader creates one deliberate Lua global, |luafemm|, to
hold the loaded module and current model. Numerical scripts can use
|local femm=require('luafemm')| without a running TeX interpreter.

\begin{description}
\item[luafemm.lua] Material laws, model ownership, assembly, Newton,
  preconditioning, field evaluation and contour topology.
\item[luafemm-boundary.lua] Side/tagged conditions, corner constraints, component gauge
  compatibility and exact P1 boundary integrals.
\item[luafemm-ideal.lua] Confined-domain compilation, floating boundary
  degrees of freedom, exact section flux and component mean contours.
\item[luafemm-integrals.lua] Segment/triangle intersection and exact integration
  of the discrete field along profiles.
\item[luafemm-shapes.lua] E and toroidal geometries, tapered poles and magnet slices.
\item[luafemm-components.lua] Pure component geometry, physical anchors and
  area-based winding currents; no TeX or solver state.
\item[luafemm-fem.lua] Passive magnetic file codec, supported-physics checks,
  source metadata and the planar model adapter; computational problem export.
\item[luafemm-topology.lua] Split planar constraints, oriented face walks,
  nested face assignment, exclusion carving and active-domain topology audits.
\item[luafemm-bh.lua] Natural cubic slopes, bounded FEMM-compatible smoothing
  and constitutive derivatives without mutating the input table.
\item[luafemm-ans.lua] Static answer-file serialization from a converged
  problem/mesh snapshot, with no remeshing or solver invocation.
\item[luafemm-cache.lua] Passive versioned disk records, exact mesh/solution
  descriptors, validation, atomic replacement and on-demand restoration.
\item[luafemm-mesh.lua] PSLG subdivision, constrained Delaunay meshing,
  local density and final mesh audit.
\item[luafemm-refine.lua] Bounded segment-first quality refinement, a triangle
  priority queue and protected circumcentre insertion.
\item[luafemm-geometry.lua] Winding/parity membership shared by refinement
  and element material classification.
\item[luafemm-predicates.lua] Filtered orientation/incircle/disk signs with
  expansion fallback; diagnostic fallback counters.
\item[luafemm-path.lua] Affine frames and soft-path polygonalisation.
\item[luafemm-materials.lua] Catalogue lookup and static interpretation;
  the separate generated data module retains upstream values.
\item[luafemm-machines.lua] Independent rotor/stator geometry, compound iron
  domains, slot validation, conductor layers and balanced winding compilation.
\item[luafemm-profiles.lua] Named measurement paths, polar and exact circular
  sampling, cached rows and discrete spatial Fourier analysis.
\item[luafemm-tex.lua] TeX boundary: escaped numeric strings, picture
  initialisation, material resolution and emitted drawing commands.
\end{description}
Public functions have LDoc-style comments with units, arguments and
return values. Private numerical helpers document invariants and algorithmic
choices rather than restating each arithmetic operation. StyLua configuration
fixes Lua 5.3 syntax, four-space indentation and a 100-column target.

\subsection{Core calls}
The public constructor copies its input; material and region declarations
also take independent copies. The structures below are owned by their model.
Callbacks and logging functions are retained as function references.
\begin{codeexample}[code only]
local femm = require('luafemm')
local model = femm.new {
  unit = .001, mesher = 'delaunay', h = 2,
  xmin = -20, xmax = 20, ymin = -20, ymax = 20,
  tolerance = 1e-7, max_newton = 60
}
femm.material(model, 'magnet', {library='cast-alnico-5-lng37'})
femm.region(model, 'magnet', {{-3,-5},{3,-5},{3,5},{-3,5}}, 0, 90)
femm.solve(model)
local bx, by, az, hx, hy = femm.sample(model, 0, 8)
\end{codeexample}

|new(options)| accepts the numerical counterparts of problem keys:
|h|, |xmin|, |xmax|, |ymin|, |ymax|, |mesher|, |mesh_tolerance|,
|min_angle|, |max_area|, |max_steiner|, |tolerance|, |max_newton|,
|depth| and |field_model| (default |open|).
Quality options use degrees, model units squared and additional vertex counts,
with defaults 0, 0 and 5000 respectively. |unit| is metres per model coordinate unit,
default .001. Optional |boundary(x,y)| and |source(x,y)| callbacks receive
SI coordinates and return $A_z$ in T\,m and $J_z$ in A/m$^2$, respectively.
The source callback replaces region currents. |log(message)| receives
iteration diagnostics. The Lua-only mesher default is |grid|.

|material(model,name,properties)| accepts |mur|, |bh|, |coercivity|,
|remanence|, |current| and |library|. |region(model,name,points,current,angle,mesh_size)|
uses model coordinates, A/m$^2$, degrees and model-unit local spacing.
Omitted current inherits the material source; omitted angle and local
spacing are zero. |mesh(model)| prepares the discrete model, and
|solve(model)| meshes if needed and returns the same solved model.

Region declarations return the owned region record. Set its |ideal_domain=true|
before meshing to include it in |field_model='ideal'|. Use
|excitation(model,x,y,NI)| for an ideal-window seed, and
|flux(model,points)| for signed right-normal section flux in Wb; the latter
requires |depth| and validates that the section stays in the active domain.
|require('luafemm-ideal').mean(model,node,cycle,kind)| returns a model-coordinate
polyline for a registered component cycle. The numerical module does not create
TeX nodes or profile names. See Section~\ref{sec:ideal} for units and limitations.

For compound regions, use
\nolinkurl{region_contours(model,name,contours,current,angle,mesh_size,fill_rule)}.
Each contour is an array of vertex pairs without a repeated closing vertex.
The fill rule is |'nonzero'| (default) or |'even odd'|. Other arguments use
the same units and defaults as |region|. Inputs are copied; every contour
must be individually simple. Model regions store their geometry in |contours|;
use |region.contours[1]| to access the first contour.
This API defines material holes, never exclusions from the PDE domain.

|sample(model,x,y)| returns |Bx,By,Az,Hx,Hy| in SI; inputs use model
coordinates. |contours(model,count,levels)| returns oriented path arrays
with |x| and |y| in model units and a |level| in T\,m. Optional explicit
|levels| replace the automatic level count. |constitutive(material,B)|
returns $\mu_0H_{\rm table}/B$ and $\mu_0\,dH_{\rm table}/dB$.
These are scaled law coefficients, not the physical vector H of a magnet.

Solved |model.stats| has the names listed in Section~\ref{sec:field}.
|model.A| is the nodal potential in T\,m, |model.nodes| contains SI positions,
and |model.triangles| contains P1 connectivity, material references and
element fields. Treat these implementation structures as read-only.
|import_mesh|, |prepare_mesh| and |update_fields| are internal preparation hooks, not a
stable foreign-mesh import interface in this release.

\subsection{Boundary data}
|boundary(model,side,condition)| replaces one side or |all| sides before
meshing. The condition table has |type| (|dirichlet|, |neumann| or |robin|),
|value|, |dx|, |dy| and |coefficient|. Omitted type means Dirichlet; omitted
numbers mean zero. The affine data are |value+dx*x+dy*y|, with $x,y$ in
metres, independently of |options.unit|. Their physical units and signs
are given in Section~\ref{sec:boundaries}. Numeric side conditions support
persistent caching. They cannot be combined with the older |boundary(x,y)|
callback, which remains a prescribed potential on the entire exterior.

The independent module |luafemm-boundary.lua| owns validation, reconstruction
of exterior edge data, gauge compatibility and consistent P1 edge assembly.
The geometric set |model.exterior_nodes| is distinct from |model.fixed_values|,
the prescribed potential map. |model.gauge_node| is present only when a
reference was needed. These prepared structures are internal and read-only.

The generic TeX layer queues |femm/boundary| options until the model's
begin-picture hook runs. A picture-local flag prevents repeated
|femm/problem| keys from replacing that model. After initialisation,
the key uses the command directly. The queue and flag follow TeX grouping;
the resulting physical declarations belong to the Lua model.

\subsection{Component geometry}
|require('luafemm-components').make(kind,options)| constructs |u core| or
|u electromagnet| geometry without a running TeX interpreter. Option names
use underscores, for example |leg_thickness|, |ampere_turns| and |gap_mesh_size|.
The result owns its options, polygon regions, named anchors and envelope
half-dimensions. It does not register a model, resolve materials or solve.
The component key reference gives the corresponding dimensions and defaults.

|current(points,unit,turns)| returns uniform current density from signed
ampere-turns and the area of a transformed polygon. The TeX bridge calls it
after capture and rounding. The PGF shape saves a geometry identifier; anchor
lookup does not depend on the current model and has no registration effects.
Only its drawing hook emits physical paths, once at the final transform.

\subsection{Persistent cache options}
Pass |cache={file='results/magnet',mode='auto'}| to |new| to enable disk
reuse. Lua filenames are relative to the process working directory; the
TeX bridge additionally honours the engine's output directory.
Omitting |cache| disables disk access. An explicit table defaults to
|mode='auto'| and requires a filename stem unless its mode is |off|.
The four modes and validation rules are described in Section~\ref{sec:cache}.

The model's |cache_info.status| is |off|, |pending|, |miss|, |mesh|,
|solution| or |refresh|. |cache_info.file| contains the resolved filename;
|cache_info.reason| gives a diagnostic when applicable. The constructor
copies cache options just like other input tables. Cache configuration
and physical model data must be treated as immutable after meshing.
Calling |mesh| or |solve| again on an already prepared model is idempotent.
Create a new model to change physical inputs or to refresh a calculation.

Persistent caching rejects |source| and |boundary| callbacks: a closure's
captured state cannot be inferred safely from its function body. These
Lua-only callbacks remain supported with |mode='off'|. Logging callbacks
are allowed and excluded from cache descriptors. Cache records contain
no executable code and never reconstruct functions.

\subsection{Catalogue and profile calls}
|require('luafemm-materials').get(id)| returns an independent record with
|id|, |name|, |path|, |fields| and |bh|. Exact original names are also accepted
when unambiguous. |list()| returns all independent records in upstream
order. |properties(id)| returns the effective static material properties.

Load the profile module with |require('luafemm-profiles')|. Its
\nolinkurl{register(model,name,points,options)} function registers a polyline with |samples|,
|offset| and |outside| options. |sample(name)| returns cached rows containing
|s,t,x,y,B,Bx,By,Bt,Bn,H,Hx,Hy,Ht,Hn,Az|. |get(name)| exposes the internal
profile record and should be treated as read-only. \nolinkurl{coordinates(name,component,abscissa)}
formats the selected columns as PGFPlots coordinate pairs. Re-register
a profile after any manual solution modification; otherwise the cached
rows deliberately remain unchanged.

\subsection{Builds, tests and source conventions}
The repository uses an installable |tex/| layout, English source comments,
English examples and documentation, and separate upstream data notices.
All build products are confined to |build/|. Development builds require
Python 3 (standard library), TeX Live tools and optionally StyLua for
format checking; installed documents need none of those development tools.

|make test| runs numerical tests, the rational predicate oracle and the
three-format integration matrix. |make examples| compiles the full examples.
|make manual| checks that all distributed examples have source/output pairs,
compiles them, crops their empty paper margins with |pdfcrop| and Ghostscript,
executes the inline examples, builds the key index and
resolves cross-references. |make materials| deterministically regenerates
the material module and catalogue from the preserved source file.
|make dist| creates local source and TDS archives; it does not publish,
commit or tag a release. The contributor guide documents the release checks.

\begin{command}{\luafemmversion}
Expandable semantic version, currently |0.6.0-dev|.
\end{command}
\begin{command}{\luafemmversiondate}
Expandable release date, currently |2026-09-27|.
\end{command}
\begin{command}{\luafemmversionnumber}
Expandable integer for simple compatibility comparisons, currently |80000|.
It increases when a release changes the public interface; it is not the
number of commits.
\end{command}

The original package code and documentation use the LaTeX Project Public
License, version 1.3c (LPPL-1.3c). Christophe Jorssen is the author and
Current Maintainer; the maintenance status is \texttt{maintained}. The
imported FEMM data retain their upstream terms. PGF's documentation
environments are loaded from TeX Live, not copied into this repository.
The manual's reference-and-example presentation was modelled on tikz-ext;
its prose and examples were written for luafemm.
